\chapter{Introducción}


\section*{Justificación de la investigación}

En las últimas tres décadas, la economía mundial se ha caracterizado por un acelerado proceso de globalización e integración regional. Esto ha favorecido la fragmentación internacional de la producción y la formación de \gls{cgv}. Como resultado de estos procesos, se ha intensificado la interdependencia entre los sectores económicos de distintos países. Ante este panorama, resulta relevante medir el nivel de interdependencia productiva y analizar la formación de bloques productivos, más allá de la existencia formal de procesos de integración económica, como las uniones aduaneras o los tratados de libre comercio.

Las tablas anuales de insumo--producto inter-países (\textit{Inter--Country Input--Output tables}, \acrshort{icio}) de la \gls{ocde} constituyen un marco contable que captura las relaciones de interdependencia productiva entre sectores económicos y países. Transformar estas tablas en redes dirigidas y ponderadas permite aplicar métodos de detección de comunidades en redes para identificar posibles bloques productivos. Bajo este enfoque, cada comunidad corresponde a un conjunto de sectores económicos, pertenecientes a uno o varios países, que mantienen fuertes vínculos productivos entre sí y vínculos más débiles con los sectores de otras comunidades. La comparación entre las comunidades identificadas y los acuerdos comerciales formales permitiría evaluar el grado de interdependencia productiva alcanzado por los procesos de integración regional, como el \gls{tlcan} o la \gls{asean}.

En esta investigación se aplican dos algoritmos de detección de comunidades para analizar las redes anuales insumo--producto internacionales construidas a partir de la base de datos \acrshort{icio}, versión 2025. Esta versión incluye información de 50 sectores económicos y 80 países para el periodo 1995--2022. En comparación con trabajos similares en la literatura, este estudio amplía la cobertura temporal, sectorial y geográfica del análisis, así como el conjunto de métodos de detección de comunidades aplicados al estudio de redes insumo--producto internacionales.

\section*{Objetivo general}

Analizar la interdependencia productiva y los procesos de integración comercial regional mediante la detección de comunidades en las redes insumo--producto internacionales durante el periodo 1995--2022.

\subsection*{Objetivos particulares}

\begin{itemize}
    \item Formalizar la representación de una matriz insumo--producto multirregional como una red dirigida y ponderada de flujos intersectoriales.
    \item Construir una serie de redes insumo--producto internacionales anuales a partir de la base de datos  \acrshort{icio} de la \gls{ocde}.
    \item Explicar cómo los métodos de detección de comunidades pueden aplicarse en redes insumo--producto para identificar bloques productivos.
    \item Identificar comunidades de sectores económicos con fuertes vínculos productivos en las redes insumo-producto internacionales entre 1995 y 2022.
    \item Analizar la estructura y evolución de las comunidades identificadas como aproximación a la interdependencia productiva y los procesos de integración comercial.
\end{itemize}
